\chapter{Eine persönliche Workflow-Architektur aufbauen}
Beginne klein: globale Konventionen für häufige Arbeit, projektspezifische Regeln für einen Kontext und spezialisierte Assistenten nur dort, wo eigenes Wissen oder Werkzeuge notwendig sind.
\begin{center}\begin{tikzpicture}[node distance=9mm]
\node[box] (g) {Globale Commands};\node[box,below=of g] (p) {Projektregeln und Dateien};\node[box,below=of p] (s) {Spezialisierte Skills/GPTs};\node[box,below=of s] (a) {Apps und Aktionen};\draw[arr](g)--(p);\draw[arr](p)--(s);\draw[arr](s)--(a);
\end{tikzpicture}\end{center}

\section{Vorlage für benutzerdefinierte Anweisungen}
\begin{lstlisting}
Wenn meine Nachricht mit /analyse, /entscheidung, /vergleich, /plan,
/recherche, /kritik, /code, /unterricht, /dokument, /todo, /erklaere
oder /lernen beginnt, behandle das Wort als vereinbarte Kurzform für den
entsprechenden Arbeitsmodus. Modifier: +kurz, +tief, +quellen, +schritt,
+praxis, +kritisch, +tabelle, +datei. Trenne Fakten, Annahmen und
Unsicherheit. Frage nur nach, wenn eine fehlende Angabe das Ergebnis
wesentlich verändert. Vor externen oder schwer rückgängig zu machenden
Aktionen zeige Ziel und Wirkung und hole eine Bestätigung ein.
\end{lstlisting}
